\documentclass[12pt]{article}
\usepackage[hidelinks]{hyperref}
\title{MINI GAME HUB REPORT}
\date{April 2026}
\author{}

\begin{document}

\maketitle
\tableofcontents
\newpage

\section{Introduction}
This project implements a secure, multi-user game hub that integrates Bash scripting for authentication and Python (Pygame) for gameplay. Two authenticated players can select a game from a menu and play via a graphical interface. The hub supports three games, Tic-Tac-Toe, Othello, and Connect Four, each rendered using Pygame with boards represented as NumPy arrays.

\section{External Libraries}
\begin{description}
    \item[numpy] Used to represent all game boards as arrays and perform win condition checks using array operations and slicing.
    \item[pygame] Used to render the game menu and all gameplay via a graphical interface window.
    \item[matplotlib] Used to display a bar chart of top players and a pie chart of most played games after each session.
    \item[pathlib] Used to read and write \texttt{history.csv} for recording game results.
    \item[sys] Used to accept command-line arguments (player usernames) passed from \texttt{main.sh} to \texttt{game.py}.
    \item[subprocess] Used to call \texttt{leaderboard.sh} from within \texttt{game.py} after each game.
    \item[datetime] Used to record the date of each game result in \texttt{history.csv}.
    \item[abc] Used to define the abstract base class for all games.
    \item[collections] Used to count frequencies for the matplotlib charts.
\end{description}

\section{Design \& Architecture}
\subsection{System structure}
The Mini Game Hub consists of two main parts: authentication and game execution.\\\\
The system starts with a Bash script (\texttt{main.sh}) which handles user login and registration. After successful authentication of two users, it runs the Python program (\texttt{game.py}).\\\\
The Python program displays the menu, allows users to select and play games, and shows the results after each game. The results are also stored in files for leaderboard purposes.\\\\
Overall flow:
\begin{itemize}
\item User authentication (Bash)
\item Game selection and execution (Python)
\item Result display and storage
\end{itemize}

\subsection{Major Components}
\begin{description}
\item[Authentication Module (main.sh)] Handles user login and registration using stored credentials.
\item[Game Engine (game.py)] Controls menu display, user interaction, game execution, and result display using Pygame.
\item[Game Modules] Each game (Tic-Tac-Toe, Othello, Connect4) is implemented separately with its own logic.
\item[Leaderboard/History] Stores and displays game results.
\item[Data Storage] 
\begin{itemize}
\item \texttt{users.tsv} for user data
\item \texttt{history.csv} for game results
\end{itemize}
\end{description}
\subsection{Design decisions}
Bash is used for authentication because it is simple for handling files and user input. Python is used for game implementation as it supports better logic handling and graphical display using Pygame.\\\\
File-based storage is used instead of a database because it is simple and sufficient for this project.\\\\
A modular design is used for games so that new games can be added easily and the system remains organized.\\\\
Passwords are stored in hashed form to improve security.
\section{Features implemented}
\begin{description}
    \item[User Authentication] Two players are asked to
 enter their username and password. Passwords are hashed using 
    SHA-256 before being stored in \texttt{users.tsv}. New users can 
    register directly from the login prompt.

    \item[Game Menu] After authentication, both players are 
    presented with a Pygame GUI menu to select from three available games.

    \item[Tic-Tac-Toe] Played on a 10×10 board where players 
    need 5 marks in a row to win, rendered using Pygame.

    \item[Othello (Reversi)] Played on an 8×8 board with 
    standard Othello rules including disc flipping and turn skipping, 
    rendered using Pygame.

    \item[Connect Four] Played on a 7×7 grid where coins 
    fall to the lowest empty row, with coins changing color every 
    move, rendered using Pygame.

    \item[Result Recording] After each game, the winner, 
    loser, date, and game name are appended to \texttt{history.csv}.

    \item[Leaderboard] \texttt{leaderboard.sh} reads 
    \texttt{history.csv} and displays a formatted table of wins, 
    losses, and win/loss ratio per player per game.

    \item[Statistics Visualisation] Matplotlib charts are 
    displayed after each game, showing the top 5 players by wins and 
    most played games.

    \item[Post-Game Loop] Players are prompted to play 
    another game or quit after viewing the statistics.
\end{description}

\section{Directory Structure}
\begin{verbatim}
mini-game-hub/
|--main.sh           #Handles user authentication and starts the game
|--game.py           #Main controller that manages game flow and user interaction.
|--leaderboard.sh    #Processes game history and displays game stats
|
|--games/
|   |--tictactoe.py  #Implements Tictactoe game logic
|   |--othello.py    #Implements othello game logic
|   |--connect4.py   #Implements Connect game logic
|       
|--users.tsv         #Stores user login credentials (hashed passwords)
|--history.csv       #Stores game results for leaderboard and analysis
\end{verbatim}

\section{Installation \& Setup}
The project requires Python 3.10, Bash, and the following Python libraries:
\begin{itemize}
    \item NumPy
    \item Pygame
    \item Matplotlib
\end{itemize}

These can be installed using:
\begin{verbatim}
pip install pygame-ce numpy matplotlib
\end{verbatim}

\begin{itemize}
    \item Ensure all project files are present in the same directory.
    \item Provide execution permissions to the scripts:
\end{itemize}

\begin{verbatim}
chmod +x main.sh
chmod +x leaderboard.sh
\end{verbatim}

\section{Usage}
Open terminal in project directory and execute
\texttt{bash main.sh}
This command initializes the system and begins the user input process.
\subsection{Login}
\begin{itemize}
    \item Upon execution the system displays welcome message.
    \item The user is prompted to enter username (and password).
    \item After successful login the system automatically proceeds to the game interface.
\end{itemize}
\subsection{Main Menu interaction}
\begin{itemize}
    \item A graphical window using pygame appears. It contains available games
    \item Clicking on the game options will launch the selected game with a small delay.
\end{itemize}
\subsection{Gameplay}
\begin{itemize}
    \item The selected game runs within the application.
    \item Players interact locally on the same system.
\end{itemize}
\subsection{Post-game}
\begin{itemize}
    \item After the game is over the same window first display the name of the winner (its a draw in case of draw)
    \item Immediately after that a leaderboard is displayed on the terminal
    \item Simultaneously, a pop-up appears and the user selects the metric to view (e.g., wins, losses, draws, ratio).
    \item Based on the selected metric, player statistics are displayed.
\end{itemize}
\subsection{Exit}
\begin{itemize}
    \item A window appears to replay or exit the game.
    \item The user can replay by clicking on play again, which displays main menu again.
    \item The user can exit the window by clicking quit and this returns you back to the terminal
\end{itemize}
\section{Retrospective}
\subsection{Hurdles faced}
The biggest challenge faced was dividing the work equally and independent of each other especially for \texttt{game.py}.\\\\
One of the hurdles faced was different styles of writing code and assigning variables. It was difficult to understand each other’s code and then debugging errors.\\\\
Another hurdle we faced was small game logic errors which we found by playing the games a lot of times trying to create all possible situations that could happen in the game.
\subsection{What would I do differently?}
If I had additional 10 hours, I would focus on refining the base class and enhancing the overall structure of the code to make it more modular and extensible. I would also work to maintain the consistency of the variable naming to make the code more readable, easier to understand, and maintain. I would also try to enhance the design and visual aspects of the game components. In addition, I would consider adding more features to the authentication system to enhance usability.\\\\
Apart from that, I would add an option of multiple themes for background, board, and pieces. This game lacks the feature of switching player turn(first player remains unchanged) after play again. So i would like to add a feature that asks if the users want to switch player turn. I would also work towards resizable pygame window.

\newpage

\begin{thebibliography}{9}
\bibitem{sha256} SHA-256 in Bash: sha256sum command
\bibitem{pygame} PYGAME CE Documentation: \url{https://pyga.me}
\bibitem{inheritance} Python Inheritance: \url{https://docs.python.org/3/tutorial/classes.html}
\bibitem{othello}Othello Rules: \url{https://worldothello.org}
\bibitem{connect4}Connect4 Rules: \url{https://boardgamegeek.com/boardgame/2719/connect-four}
\bibitem{numpy}Numpy Documentation: \url{https://numpy.org/doc/2.4/}
\end{thebibliography}

\end{document}
